%% Observation Theory for Estimation and Control -- first LaTeX draft.
%% Supersedes paper/Observation Theory for Estimation and Control.docx,
%% implementing paper/OT-EC-REVISION-NOTES.md (v2, 2026-08-18) and the
%% prior-art sweep paper/PRIOR-ART-SWEEP-estimation-control.md.
\documentclass[letterpaper,onecolumn]{IEEEtran}

\usepackage[T1]{fontenc}
\usepackage[utf8]{inputenc}
\usepackage{amsmath,amssymb,amsthm,mathtools}
\usepackage{booktabs}
\usepackage{array}
\usepackage{bm}
\usepackage{microtype}
\usepackage{cite}
\usepackage{xcolor}
\usepackage[colorlinks=true,linkcolor=blue,citecolor=blue,urlcolor=blue]{hyperref}
\hypersetup{pdftitle={Observation Theory for Estimation and Control},pdfauthor={Andrew H. Bond}}

\newtheorem{theorem}{Theorem}
\newtheorem{proposition}{Proposition}
\newtheorem{lemma}{Lemma}
\newtheorem{corollary}{Corollary}
\theoremstyle{definition}
\newtheorem{definition}{Definition}
\newtheorem{conjecture}{Conjecture}
\theoremstyle{remark}
\newtheorem{remark}{Remark}

\newcommand{\E}{\mathbb{E}}
\newcommand{\R}{\mathbb{R}}
\newcommand{\PC}{P_{\mathcal{C}}}
\newcommand{\PCb}{\bar P_{\mathcal{C}}}
\newcommand{\tr}{\operatorname{tr}}
\newcommand{\C}{\mathcal{C}}
\newcommand{\Sig}{\Sigma}
\newcommand{\Sm}{\Sigma^{-}}
\newcommand{\Sp}{\Sigma^{+}}
\newcommand{\VC}{V_{\mathcal{C}}}
\newcommand{\UC}{U_{\mathcal{C}}}
\newcommand{\SC}{S_{\mathcal{C}}}

\begin{document}

\title{Observation Theory for Estimation and Control:\\
Consumer-Induced Geometry as a Complement to\\
Observability, Filtering, and Feedback}

\author{Andrew~H.~Bond,~\IEEEmembership{Senior~Member,~IEEE}%
\thanks{First LaTeX draft, August 2026. The author is with the Department of
Computer Engineering, San Jos\'e State University, San Jos\'e, CA 95192 USA
(e-mail: andrew.bond@sjsu.edu; ORCID 0009-0003-2599-6158).}%
\thanks{This manuscript extends the author's Observation Theory program from
compression and coding to estimation and control. Program repository, sealed
preregistrations, claims ledger, and the Lean~4/Mathlib formalization of
Proposition~\ref{prop:invariance} and Proposition~\ref{prop:voi} are at
\url{https://github.com/ahb-sjsu/geometric-observation}.}}

\maketitle

\begin{abstract}
Classical control theory separates several questions that are easy to
conflate: whether a plant can be driven, whether its state can be inferred,
how uncertainty should be filtered, and how actions should be chosen. A
further question is usually supplied externally by the designer: \emph{which
distinctions in the estimated state matter to the downstream consumer?}
Observation Theory (OT) supplies a candidate mathematical object for that
missing interface. If a consumer $\C:\R^d\to\R^m$ is judged in a local output
metric $G$, its pullback $\PC(x)=J_\C(x)^{\top}G(\C(x))\,J_\C(x)$ induces the
local state geometry the consumer can distinguish, and estimator uncertainty
$\Sig$ and consumer geometry compose as the operational uncertainty
$\UC=\tr(\PC\Sig)$. This paper defines a research campaign around that
composition, organized as three interfaces to the classical stack:
\emph{estimator geometry} (a state-dependent $\PC(x)$ moves the optimal
estimate off the posterior mean), \emph{resource value}
($\VC=\tr(\PC(\Sm-\Sp))$, the consumer-relative value of a measurement), and
\emph{dynamic relevance} ($\SC(\Delta)=\tr[\PC\,\Sig(t{+}\Delta\mid I_t)]$,
the age of operationally relevant uncertainty). We first establish a
boundary: for a fixed positive-semidefinite $\PC$ and an unconstrained
Bayesian estimator, the posterior mean --- hence the Kalman estimate ---
remains optimal; both this invariance and the rank-one value-of-observation
identity are machine-checked in Lean~4/Mathlib. Each pillar of the program
has close prior art in isolation: black-box recovery of quadratic task
metrics (decision-focused learning), and consumer-derived weights composed
with covariances (LQG sensing co-design, goal-oriented experimental design,
value-of-information scheduling) in analytic form. The narrow novelty
hypothesis is their conjunction: that a downstream consumer's operational
metric can be recovered from the consumer itself under \emph{query access
only}, composed with classical estimator uncertainty, and validated
prospectively, out of sample, at matched budgets that include the cost of
probing. OT is proposed not as a replacement for observability, filtering,
stability, or feedback, but as a complementary theory of operational
distinguishability at their downstream boundary.
\end{abstract}

\begin{IEEEkeywords}
Kalman filtering, sensor scheduling, value of information, task-oriented
communication, goal-oriented experimental design, functional observability,
operational distinguishability.
\end{IEEEkeywords}

\section{The missing question in the classical control stack}
\label{sec:intro}

A modern estimation-and-control stack can be schematized as
\begin{equation*}
\text{plant} \;\to\; \text{sensors} \;\to\; \text{estimator} \;\to\;
\text{consumer(s)} \;\to\; \text{actions}.
\end{equation*}
Classical systems theory provides precise concepts for controllability and
observability. Kalman filtering \cite{kalman1960} provides an optimal state
estimate and error covariance under its linear-Gaussian assumptions. LQR,
LQG, MPC, and stochastic control determine actions once an objective has been
specified. Yet a design choice remains embedded in these methods: \emph{which
state errors actually matter?} In LQR this appears as the choice of a state
penalty $Q$; in estimation it appears in accuracy specifications, sensor
priorities, covariance weights, update rates, and communication policies. In
many systems these choices are engineered successfully --- but their
operational geometry is externally specified.

Observation Theory begins one step downstream. Let $\C:\R^d\to\R^m$ be a
consumer and let $D_Y$ be a twice-differentiable divergence on its output.
Define the output metric and consumer Jacobian
\begin{equation}
G(z)\;=\;\nabla_v^2\,D_Y(z,v)\big|_{v=z},
\qquad
J_\C(x)\;=\;\frac{\partial \C}{\partial x}(x),
\end{equation}
and the \emph{read operator}
\begin{equation}
\label{eq:pullback}
\boxed{\;\PC(x)\;=\;J_\C(x)^{\top}\,G(\C(x))\,J_\C(x).\;}
\end{equation}
For sufficiently small state error $\delta=x-\hat x$,
\begin{equation}
D_Y\!\big(\C(x),\C(\hat x)\big)\;\approx\;\tfrac12\,\delta^{\top}\PC(x)\,\delta,
\qquad
\E\,D_Y \;\approx\; \tfrac12\,\tr\!\big(\PC\,\Sig_\delta\big),
\end{equation}
so ordinary reconstruction error is the isotropic corner $\PC=I$. The
composition
\begin{equation}
\label{eq:composition}
\UC \;=\; \tr\!\big(\PC\,\Sig\big)
\end{equation}
reads: \emph{what the estimator does not know}, weighted by \emph{what the
consumer can distinguish}. Neither object substitutes for the other: an
accurate consumer model cannot estimate the plant; a perfect error covariance
cannot say which residual uncertainty matters operationally.

The systems question is therefore not ``can Observation Theory replace the
Kalman filter?'' It is: \emph{given a classical estimator that tells us what
we do not know, can a consumer-induced geometry tell us which parts of that
uncertainty matter --- and therefore where scarce sensing, communication,
computation, and control resources should be spent?}

\subsection{Three interfaces, not one}
\label{sec:interfaces}

The program is organized around three distinct interfaces between the read
operator and the classical stack:
\begin{equation}
\label{eq:threeinterfaces}
\PC(x)\;\longrightarrow\;
\begin{cases}
\textbf{Estimator geometry} & \text{if $\PC$ varies with state
  (\S\ref{sec:boundary}, Rem.~\ref{rem:statedep})},\\[2pt]
\textbf{Resource value} & \tr\!\big(\PC\,\Delta\Sig\big)
  \quad\text{(\S\ref{sec:voi}--\S\ref{sec:resources})},\\[2pt]
\textbf{Dynamic relevance} & \tr\!\big(\PC\,\Sig(\Delta)\big)
  \quad\text{(\S\ref{sec:staleness}, conjectural)}.
\end{cases}
\end{equation}
Interface~1 is \emph{opened} rather than closed by the boundary theorem of
\S\ref{sec:boundary}: the posterior-mean invariance holds for fixed $\PC$,
and its failure under state-dependent $\PC(x)$ is precisely where consumer
geometry can change the estimator itself. Interface~2 is the paper's spine.
Interface~3 is developed as an explicitly conjectural branch.

\subsection{Access models: demonstrations, gradients, queries}
\label{sec:access}

The distinction that defines this program's niche is the access model.
Inverse optimal control and inverse reinforcement learning recover objectives
from \emph{demonstrations}. Decision-focused learning differentiates through
a \emph{known} (or differentiable) consumer
\cite{elmachtoub2022,amos2017,mandi2024}. Observation Theory probes a
\emph{black-box} consumer under a declared observation distribution. Only the
third access model is available when the consumer is a learned perception
stack, a physical link, or another party's system --- and those consumers,
which resist hand-modeling, are exactly where a recovered geometry is
necessary rather than convenient.

\section{Observability is not operational distinguishability}
\label{sec:observability}

Consider $x_{k+1}=Ax_k+Bu_k+w_k$, $y_k=Cx_k+v_k$. Classical observability
asks whether distinct internal states can be distinguished from the
measurement history. OT asks a downstream question: for consumer $\C$, two
nearby states $x$ and $x+\delta$ are \emph{locally operationally
distinguishable} according to $\delta^{\top}\PC(x)\,\delta$. A direction in
$\ker\PC$ is locally invisible to that consumer; a direction with a large
eigenvalue of $\PC$ is strongly read. Hence a state direction can be
perfectly observable to the sensor system and nearly irrelevant to the
consumer; conversely, a very small estimation error can be operationally
dominant if it lies in a strongly weighted read direction. These are
complementary geometries on opposite sides of the estimator.

Two classical structures must be engaged here by name. First,
\emph{functional observability} \cite{darouach2000,montanari2022} asks when a
hand-specified linear functional $z=Lx$ can be reconstructed; OT's linear
special case $\C(x)=Lx$ lies inside its subject matter, and OT adds the
\emph{metric} (not merely subspace) structure, nonlinear consumers, and
recovery of $L$-like structure from the consumer itself. Second, for a linear
consumer $z=Lx$ with output metric $G$, the pointwise OT geometry is
$\PC=L^{\top}GL$, and the output-weighted observability Gramian of
frequency-weighted balanced truncation \cite{enns1984} is its
\emph{dynamics-propagated accumulated analogue},
\begin{equation}
\label{eq:gramian}
W_\C(T)\;=\;\int_0^{T} e^{A^{\top}t}\,\PC\,e^{At}\,dt .
\end{equation}
Classical observability geometry thus \emph{emerges from accumulated OT
geometry} in the linear case: the Gramian is what the consumer's
instantaneous read metric integrates to under the flow.

\section{A boundary proposition: OT does not create a new Kalman filter}
\label{sec:boundary}

The first result of the campaign is deliberately negative.

\begin{proposition}[Weighted posterior-mean invariance]
\label{prop:invariance}
Let $X$ be a random state, $Y$ the available observations, and $P\succeq 0$ a
fixed state weighting. Then
$\E\big[(X-\hat x)^{\top}P\,(X-\hat x)\mid Y\big]$
is minimized over $\hat x$ by the posterior mean $\hat x=\E[X\mid Y]$.
\end{proposition}

\begin{proof}
Let $m=\E[X\mid Y]$. Then $X-\hat x=(X-m)+(m-\hat x)$, the cross terms vanish
in conditional expectation, and
\begin{equation*}
\E\big[(X-\hat x)^{\top}P(X-\hat x)\mid Y\big]
=\tr\!\big(P\,\Sig_{X\mid Y}\big)+(m-\hat x)^{\top}P\,(m-\hat x),
\end{equation*}
whose second term is nonnegative.
\end{proof}

The posterior mean already minimizes \emph{every} fixed positive-semidefinite
quadratic state loss; in the linear-Gaussian case this is the Kalman
estimate. If the campaign merely rediscovered a weighted Kalman filter, it
would not constitute the proposed contribution. The opportunity lies in
deciding which measurements to take, which information to transmit, which
state directions to retain, when to update, and eventually which operational
errors to regulate --- and, per Remark~\ref{rem:statedep}, in the
state-dependent regime the invariance itself fails.

\begin{remark}[Machine-checked form; weakened hypotheses]
\label{rem:kernel}
Proposition~\ref{prop:invariance} (discrete form) and
Proposition~\ref{prop:voi} below are machine-checked in Lean~4 against
Mathlib (\texttt{lean/ObservationTheory/WeightedMeanInvariance.lean}; Lean
and Mathlib v4.32.2, clean build, no \texttt{sorry}). The formalization shows
the invariance survives under hypotheses weaker than stated: only
$v^{\top}Pv\ge 0$ for all $v$ and weights summing to one are used --- neither
symmetry of $P$ nor nonnegativity of the weights. We state the natural
$P\succeq 0$ form above; the weakening is recorded here because the kernel
checked it, not because the generality is operationally needed. Per the
program's charter, mechanical verification is necessary but not sufficient
for a \texttt{[proved]} ledger classification, which additionally requires an
independent fresh-context derivation pass.
\end{remark}

\begin{remark}[State dependence opens Interface 1]
\label{rem:statedep}
Proposition~\ref{prop:invariance} is for \emph{fixed} $P$. For a nonlinear
consumer, $\PC(x)$ depends on $x$; the loss
$\E[(X-\hat x)^{\top}\PC(X)(X-\hat x)\mid Y]$ is no longer a fixed quadratic,
and its minimizer is in general \emph{not} the posterior mean. Consumer
geometry can then change the estimator itself, not only the allocation of
resources around it. Characterizing this estimator shift is an open problem
of the campaign, not a result of this paper.
\end{remark}

\section{Consumer-relative value of observation}
\label{sec:voi}

Let a Bayesian or Kalman estimator hold prior covariance $\Sm$ and consider a
candidate measurement $y=Hx+v$, $v\sim\mathcal N(0,R)$. The covariance update
gives $\Sp=\Sm-\Sm H^{\top}(H\Sm H^{\top}+R)^{-1}H\Sm$. Define the
\emph{consumer-relative value of observation}
\begin{equation}
\label{eq:voi}
\VC \;=\; \tr\!\big(\PC\,(\Sm-\Sp)\big)
\;=\; \tr\!\big(\PC\,\Sm H^{\top}(H\Sm H^{\top}+R)^{-1}H\Sm\big).
\end{equation}

\begin{proposition}[Rank-one form]
\label{prop:voi}
For a scalar measurement $y=h^{\top}x+v$ with $\operatorname{Var}(v)=r$ and a
rank-one consumer $\PC=gg^{\top}$,
\begin{equation}
\label{eq:voirankone}
\VC \;=\; \frac{\big(g^{\top}\Sm h\big)^{2}}{\,h^{\top}\Sm h+r\,}.
\end{equation}
\end{proposition}

This form is the section's message: a sensor is not operationally valuable
because it observes a high-variance direction; it is valuable when its
measurement direction $h$, \emph{propagated through the current uncertainty}
$\Sm$, overlaps a direction $g$ the consumer actually reads. The covariance
algebra is classical. The proposed OT contribution is the \emph{provenance}
of $g$ or $\PC$: supplied by the downstream consumer's distinguishability
geometry rather than chosen as an engineering weight.

\section{Resource-constrained operation is the natural interface}
\label{sec:resources}

Weighted covariance objectives are established in sensor selection and
scheduling \cite{joshiboyd2009,unified2023} and in LQG sensing co-design
\cite{tzoumas2021}; the optimization $\min_a \tr(M\Sp(a))$ for a
designer-specified $M$ is not an OT novelty. The proposed pipeline is
\begin{equation*}
\text{consumer} \;\xrightarrow{\;\text{probe}\;}\; \widehat{\PC}
\;\xrightarrow{\;\text{compose}\;}\; \tr\big(\widehat{\PC}\,\Sig\big)
\;\xrightarrow{\;\text{allocate}\;}\; a^{\star},
\end{equation*}
with the allocation judged on the actual downstream consumer.

\subsection{Probe cost is part of the objective}
\label{sec:probecost}

Recovering $\widehat{\PC}$ --- and re-recovering it along trajectories when
the consumer is nonlinear --- consumes exactly the resources OT allocates.
The policy problem is therefore not
$\min_\pi D_\C(\pi)$ subject to side constraints, but
\begin{equation}
\label{eq:policy}
\min_{\pi}\;\; D_\C(\pi)
\;+\;\lambda_s\,C_{\mathrm{sensing}}(\pi)
\;+\;\lambda_c\,C_{\mathrm{comm}}(\pi)
\;+\;\lambda_p\,C_{\mathrm{probing}}(\pi),
\end{equation}
which contains a genuine exploration problem: \emph{when is another consumer
query worth its cost?} The marginal value of a probe --- how much a further
refinement of $\widehat{\PC}$ changes downstream allocations --- is itself an
OT quantity. All campaign comparisons in \S\ref{sec:campaigns} are run under
\eqref{eq:policy}, with probe cost on the same budget line as sensing and
communication.

\subsection{Greedy selection: a stated limitation}
\label{sec:greedy}

Weighted-trace objectives $\tr(\PC\,\Delta\Sig)$ are \emph{not} supermodular
in general: the counterexamples of Jawaid and Smith \cite{jawaidsmith2015}
cover trace-of-covariance objectives, so greedy selection under
\eqref{eq:voi} does not inherit the classical $(1-1/e)$ guarantee that
log-determinant objectives enjoy. We do not attempt to repair this here.
Available machinery includes approximate-supermodularity bounds
\cite{chamon2021} and convex relaxations \cite{joshiboyd2009}; establishing
\emph{conditions under which $\VC$-greedy selection retains approximation
guarantees} is registered as a campaign theorem --- a named target, not a
claimed result.

\subsection{Operational event triggering}

An operational trigger updates when the uncertainty visible to the consumer
becomes excessive --- $\tr(\PC\,\Sig_t)>\tau$ --- rather than when a timer
expires or isotropic variance crosses a threshold. The analytic special case
is the value-of-information trigger of Soleymani, Baras, and Hirche
\cite{soleymani2022,soleymani2023}, in which the transmit rule derived from
the downstream LQG cost is provably optimal. Any OT triggering claim must
cite that line as its analytic ancestor and either match it (LQG) or
generalize beyond it (non-analytic consumers).

\section{Consumer geometry as a local control geometry}
\label{sec:control}

If the ultimate objective is expressed at a downstream output $z=\C(x)$ with
desired value $z^{\star}=\C(x^{\star})$, a second-order expansion gives
$D_Y(\C(x),\C(x^{\star}))\approx\tfrac12 (x-x^{\star})^{\top}
\PC(x^{\star})(x-x^{\star})$, so $\PC(x^{\star})$ is a natural local state
penalty --- a candidate \emph{provenance} for one component of an LQR/MPC
objective, not a replacement for it. For a nonlinear consumer the penalty is
trajectory dependent. Stability, robustness, terminal costs, actuator
constraints, and feasibility remain control-theoretic problems; the proposed
division of labor is that OT supplies where the state penalty comes from, and
control theory determines what to do with it.

\subsection{Threshold consumers and the belief-averaged read operator}
\label{sec:belief}

The local metric \eqref{eq:pullback} presumes a Hessian. Threshold-like
consumers --- bit-error rate, outage, any cliff --- have degenerate or
explosive $G$ near the threshold, and the paper's own application exemplar
(\S\ref{sec:fso}) is of this type. Where the local geometry exists, the
robustified object is the \emph{belief-averaged read operator}
\begin{equation}
\label{eq:beliefavg}
\PCb(b)\;=\;\E_{X\sim b}\!\left[\,J_\C(X)^{\top}\,G(\C(X))\,J_\C(X)\,\right],
\end{equation}
where $b$ is the estimator's belief: the read operator is averaged over the
uncertainty the estimator actually has, rather than evaluated at a point the
system may not occupy. For genuinely nonsmooth outage metrics no Hessian
exists to average, and a finite-perturbation or smoothed operational metric
must replace \eqref{eq:pullback} rather than pretending one does.
\emph{Belief-relative Observation Theory} --- the read operator as a
functional of the belief --- is flagged as a substantial extension in its own
right.

\section{Black-box recovery is the sharp claim}
\label{sec:recovery}

If $\PC$ must always be hand-specified, much of this program reduces to
established task-weighted estimation and control. The stronger proposition is
that the read geometry can, under stated conditions, be identified from the
consumer itself: for appropriate consumers, finite-difference probes under a
declared observation distribution recover sensitivity information without
privileged access to the consumer's implementation. The probing
\emph{methodology} is itself prior art --- empirical observability Gramians
recover plant sensitivity by perturb-and-simulate \cite{krenderide2009} ---
and in machine learning, locally optimized decision losses recover quadratic
task-metric surrogates from a black-box decision oracle by sampled
perturbations \cite{shah2022,shah2024}. What is claimed here is the probed
\emph{object} (a downstream consumer's operational metric), its
\emph{composition} with an explicit estimator covariance, and the
\emph{validation protocol}: recovery on one perturbation distribution
followed by success only on the same perturbations is insufficient; the test
must be out of sample, at matched budgets, on the held-out consumer.

\section{Multiple consumers}
\label{sec:multi}

A real state service rarely has one downstream consumer. Given
$P_1,\dots,P_m$ for navigation, tracking, collision avoidance,
communications, or mission logic, a weighted sum $\PC^{\mathrm{tot}}=\sum_i
\lambda_i P_i$ can be useful but can hide incompatibility. Observation
Theory's two-observer rate--distortion result supplies a precedent: in the
Gaussian weighted-quadratic problem, successive refinement is lossless
exactly when the observers' optimal residual uncertainty ellipsoids nest;
distinct observers can demand geometrically incompatible information. That
result is a coding theorem, not yet a control theorem. The corresponding
systems question --- when can one estimator, sensor schedule, or state
representation efficiently serve several consumers, and when does their
operational geometry force a multi-consumer tax? --- remains to be
discovered, and must not be answered by smuggling the coding result across.

\section{Dynamic relevance: the age of operationally relevant uncertainty}
\label{sec:staleness}

Interface 3 is stated as an explicit conjecture. Let $\Sig(t{+}\Delta\mid
I_t)$ be the uncertainty predicted from currently retained information.
Define
\begin{equation}
\label{eq:staleness}
\SC(\Delta)\;=\;\tr\!\big[\,\PC(t{+}\Delta)\;\Sig(t{+}\Delta\mid I_t)\,\big].
\end{equation}
An observation becomes stale when the uncertainty it leaves \emph{in the
consumer's read directions} becomes operationally expensive --- age of
information \cite{sun2020,maatouk2022} weighted by operational geometry
rather than by the clock. Our prior-art sweep found no anisotropic-staleness
formulation; the nearest neighbor is the analytic (scalar-LQG)
value-of-information line \cite{soleymani2022}. Because \eqref{eq:staleness}
has its own community, baselines (AoI, AoII, VoI), and falsifier (does
directional staleness beat age-based scheduling at matched update budgets on
a real consumer?), it is retained here only as a conjectural branch and
pursued independently. If long-horizon consequences dominate immediate
consumer-relative covariance reduction, $\SC$ may remain a feature or
terminal metric rather than a sufficient scheduling policy; that outcome
would be a scientific boundary, not a failure to be patched.

\section{Prior art and the novelty hypothesis}
\label{sec:priorart}

The foundational claim must be narrower than several established areas, and
--- following a dedicated three-slice sweep (control theory; machine
learning; communications and the application domain) --- narrower still than
the original draft of this paper stated. Each pillar of the hypothesis exists
separately in strong form. Local quadratic task metrics are recovered from
black-box decision oracles by sampled probing in decision-focused learning
\cite{shah2022,shah2024,bansal2023}. Consumer-derived weights composed with
estimator covariances drive prospective sensor and bit allocation in LQG
sensing co-design \cite{tzoumas2021}, goal-oriented Bayesian experimental
design \cite{attia2018,wu2023,quadgoed2025} --- including Jacobian/Hessian
pullbacks of nonlinear goal maps with adjoint access --- value-of-information
scheduling \cite{soleymani2022,soleymani2023,soleymani2024}, rate-cost
control \cite{kostina2019,tanaka2015}, goal-oriented quantization
\cite{zou2022}, and task-based quantization \cite{shlezinger2019}. Attention
and anticipation select sensing resources against task-projected covariance
metrics in closed loop \cite{carlone2018}; task-driven representations are
co-designed by information bottlenecks \cite{pacelli2019}; consumer-derived
perceptual metrics already steer codec bit allocation in practice
\cite{zhang2018}. No novelty is claimed for any of these, nor for Kalman
filtering, posterior-mean optimality, PSD-weighted covariance objectives,
sensor selection and scheduling, active sensing, semantic communication,
loss-calibrated inference \cite{lacostejulien2011}, or inverse optimal
control. The analytic LQG case is conceded entirely.

\begin{quote}
\textbf{Research hypothesis (conjunction).} A downstream consumer induces an
operational state geometry $\PC=J^{\top}GJ$ that can, under explicit
observation conditions, be recovered from the consumer itself under
\emph{query access only} and a declared observation distribution; composed
with classical estimator uncertainty as $\tr(\PC\Sig)$; and used to predict
and optimize the operational value of sensing, communication,
representation, and local control resources for consumers \emph{outside the
analytic families} --- validated prospectively, out of sample, at matched
budgets that include the cost of probing, on the actual downstream consumer.
\end{quote}

This is a novelty hypothesis, not a novelty verdict: no single work found
combines black-box probed consumer geometry, explicit covariance
composition, and prospective matched-budget validation, but a referee
holding decision-focused learning in one hand and goal-oriented experimental
design in the other could reasonably call the conjunction an integration.
The answer is the integration \emph{plus the validation protocol} ---
reconstruction-matched flips, held-out consumers, matched budgets, physical
endpoints --- which no found work runs.

\section{Research campaign}
\label{sec:campaigns}

\subsubsection*{Campaign 1 --- Verify the boundary}
Random linear systems; direct optimization over admissible linear estimator
gains compared with the Kalman solution under randomly drawn PSD consumer
metrics. Any repeatable improvement over the Kalman estimate indicates an
implementation error or a model violation. The invariance is classical and
now machine-checked (Remark~\ref{rem:kernel}); the campaign \emph{verifies}
the boundary publicly rather than discovers it.

\subsubsection*{Campaign 2 --- Consumer-relative value of observation}
A dynamical state with heterogeneous candidate sensors and a planted low-rank
consumer; selection by \eqref{eq:voi}. The signature experiment is a
\emph{reconstruction-matched flip}: two sensing policies with essentially
identical $\tr\Sig$ but differently shaped uncertainty; only the consumer
changes. Preregistered prediction: $\tr(\PC\Sig)$ identifies which policy is
better for the downstream consumer while ordinary state MSE cannot. Where an
analytic model exists, baselines must include the Riccati/VoI-derived weight
--- not merely hand-tuned or isotropic alternatives.

\subsubsection*{Campaign 3 --- Blind recovery, then scheduling (flagship)}
The point at which the pieces stop resembling analytic LQG/VoI machinery:
\begin{equation*}
\text{unknown consumer}\;\to\;\text{query-only recovery of }\widehat{\PC}
\;\to\;\text{sensor scheduling}\;\to\;\text{held-out real consumer},
\end{equation*}
at equal sensing budgets with probe cost included per \eqref{eq:policy}, and
with an \emph{analytic-consumer positive-control arm}: on a consumer whose
optimal weight is known (LQG/VoI), the blind pipeline is required to
\emph{recover and match} the known optimum --- not beat it. Matching the
analytic optimum blind is the calibration proof; winning on non-analytic
consumers is the contribution. Success requires more than subspace recovery:
the recovered operator must choose useful measurements and improve the
held-out consumer at matched budget.

\subsubsection*{Campaign 4 --- Operational event triggering}
Limited transmissions or measurement energy; compare periodic, age-based,
isotropic-covariance, hand-weighted, analytic-VoI (where available), and
recovered-OT triggering. Primary outcome: downstream consumer performance at
matched communication or energy --- never lower OT surrogate loss by
construction.

\subsubsection*{Campaign 5 --- Physical consumer}
A moving RF or free-space optical tracking system (\S\ref{sec:fso}).
Navigation state includes position, velocity, attitude, angular rate, and
structural motion; the physical consumer reads angle error, optical coupling,
received power, BER, or outage. The key test is again reconstruction-matched:
two navigation systems with similar state RMSE but different covariance
geometry should produce different physical performance, and a blindly
recovered $\PCb$ should predict the ordering \emph{before} the run.

\subsubsection*{Campaign 6 --- Multi-consumer state service}
Two consumers with read subspaces ranging from aligned to orthogonal; compare
common estimation, weighted scalarization, separate service, and
layered/refined descriptions. The two-observer coding theorem supplies
hypotheses about incompatibility but must not be imported as a control
result; the goal is to determine whether a genuinely dynamic analogue of
observer complementarity exists.

\subsubsection*{Campaign 7 --- Closed-loop control}
Only after the estimation and sensing claims survive: a recovered or derived
$\PC(x)$ inside an LQR/MPC objective against isotropic and competently
hand-tuned baselines at matched control effort, with stability and
constraint satisfaction independently audited, and the endpoint again the
real downstream consumer.

\section{Falsifiers}
\label{sec:falsifiers}

The program narrows or fails if: the recovered read operator does not
generalize outside its probing distribution; a conventional
engineer-specified task matrix performs as well, making recovery
operationally unnecessary; $\PC(x)$ varies too rapidly for a local metric to
hold over realistic estimator uncertainty (the belief-averaged operator
\eqref{eq:beliefavg} is the first line of defense, and its failure would
bound belief-relative OT as well); probe cost under \eqref{eq:policy}
consumes the advantage the recovered geometry provides; or long-horizon
closed-loop consequences systematically defeat policies based on
consumer-relative covariance reduction, in which case OT remains a theory of
consumer-relative representation with local and terminal uses. These
outcomes are scientific boundaries, to be reported, not patched away.

\section{Application exemplar: cooperative optical high-altitude platforms}
\label{sec:fso}

A useful application --- not the foundation of the theory --- is a
constellation of persistent high-altitude platforms connected by free-space
optical links \cite{kaymak2018}. Each vehicle has a navigation and attitude
state; each optical terminal reads only particular combinations of
transverse displacement, line-of-sight angle, attitude, angular velocity,
structural jitter, range, and optical state. Its physical consumer geometry
need not resemble 6-DOF navigation MSE, and the standard practice in the
pointing-acquisition-tracking literature is a static top-down pointing-error
budget: the causal arrow runs forward (attitude error $\to$ link budget)
only. The OT proposal runs it backward: probe the black-box receiver through
modeled state perturbations, recover a belief-averaged read operator
\eqref{eq:beliefavg} --- the endpoint metrics are threshold-like, so the
smoothed construction of \S\ref{sec:belief} is required, not optional ---
and contract it with the platform's Kalman covariance to decide when to
request a beacon update, which navigation measurement to acquire, which
information to transmit to a neighbor, and how much precision to allocate to
each state direction. Multiple simultaneous links provide a natural
multi-consumer test. The advantage of this exemplar is that its endpoint ---
received power, fiber coupling, BER, outage --- is external and physical: OT
cannot win merely by optimizing its own metric.

\section{Conclusion}
\label{sec:conclusion}

Observation Theory proposes an estimation-and-control research program
centered on a compact composition, $\UC=\tr(\PC\Sig)$: the covariance
belongs to estimation theory, the read operator to the consumer's
operational geometry; the former says what is uncertain, the latter what
matters. Three interfaces organize the program --- estimator geometry when
$\PC$ varies with state, resource value through $\tr(\PC\,\Delta\Sig)$, and
dynamic relevance through $\SC(\Delta)$ --- and none replaces observability,
filtering, LQG, MPC, or stability theory. The campaign succeeds only if its
sharpest claim survives: that $\PC$ can be recovered from real downstream
consumers under query access, at accounted probe cost, and prospectively
predicts resource allocations that ordinary state reconstruction does not
--- matching the analytic optimum where one exists and winning where none
does. The program is judged by reconstruction-matched flips, blind recovery,
held-out dynamics, physical endpoints, competent classical baselines, and
explicit negative results. If those gates fail, Observation Theory remains a
theory of consumer-relative representation. If they pass, operational
distinguishability becomes a working complement to the classical control
stack.

\begin{thebibliography}{99}\footnotesize

\bibitem{kalman1960} R.~E. Kalman, ``A new approach to linear filtering and
prediction problems,'' \emph{J. Basic Eng.}, vol.~82, no.~1, pp.~35--45, 1960.

\bibitem{darouach2000} M.~Darouach, ``Existence and design of functional
observers for linear systems,'' \emph{IEEE Trans. Autom. Control}, vol.~45,
no.~5, pp.~940--943, 2000.

\bibitem{montanari2022} A.~N. Montanari, C.~Duan, L.~A. Aguirre, and A.~E.
Motter, ``Functional observability and target state estimation in large-scale
networks,'' \emph{Proc. Natl. Acad. Sci. USA}, vol.~119, no.~1, e2113750119,
2022.

\bibitem{enns1984} D.~F. Enns, ``Model reduction with balanced realizations:
An error bound and a frequency weighted generalization,'' in \emph{Proc. IEEE
Conf. Decision and Control}, 1984, pp.~127--132.

\bibitem{krenderide2009} A.~J. Krener and K.~Ide, ``Measures of
unobservability,'' in \emph{Proc. IEEE Conf. Decision and Control}, 2009,
pp.~6401--6406.

\bibitem{tzoumas2021} V.~Tzoumas, L.~Carlone, G.~J. Pappas, and A.~Karaman,
``LQG control and sensing co-design,'' \emph{IEEE Trans. Autom. Control},
vol.~66, no.~4, pp.~1468--1483, 2021.

\bibitem{carlone2018} L.~Carlone and S.~Karaman, ``Attention and anticipation
in fast visual-inertial navigation,'' \emph{IEEE Trans. Robotics}, vol.~35,
no.~1, pp.~1--20, 2018.

\bibitem{pacelli2019} V.~Pacelli and A.~Majumdar, ``Task-driven estimation and
control via information bottlenecks,'' in \emph{Proc. IEEE Int. Conf.
Robotics and Automation}, 2019, pp.~2061--2067.

\bibitem{joshiboyd2009} S.~Joshi and S.~Boyd, ``Sensor selection via convex
optimization,'' \emph{IEEE Trans. Signal Process.}, vol.~57, no.~2,
pp.~451--462, 2009.

\bibitem{jawaidsmith2015} S.~T. Jawaid and S.~L. Smith, ``Submodularity and
greedy algorithms in sensor scheduling for linear dynamical systems,''
\emph{Automatica}, vol.~61, pp.~282--288, 2015.

\bibitem{chamon2021} L.~F.~O. Chamon, A.~Amice, and A.~Ribeiro,
``Approximately supermodular scheduling subject to matroid constraints,''
\emph{IEEE Trans. Autom. Control}, vol.~67, no.~3, pp.~1384--1396, 2022.

\bibitem{unified2023} L.~Ye \emph{et al.}, ``A unified approach to optimally
solving sensor scheduling and sensor selection problems in Kalman
filtering,'' arXiv:2304.02692, 2023.

\bibitem{soleymani2022} T.~Soleymani, J.~S. Baras, and S.~Hirche, ``Value of
information in feedback control: Quantification,'' \emph{IEEE Trans. Autom.
Control}, vol.~67, no.~7, pp.~3730--3737, 2022.

\bibitem{soleymani2023} T.~Soleymani, J.~S. Baras, S.~Hirche, and K.~H.
Johansson, ``Value of information in feedback control: Global optimality,''
\emph{IEEE Trans. Autom. Control}, vol.~68, no.~6, pp.~3641--3647, 2023.

\bibitem{soleymani2024} T.~Soleymani, J.~S. Baras, and K.~H. Johansson,
``Foundations of value of information: A semantic metric for networked
control systems tasks,'' arXiv:2403.11927, 2024.

\bibitem{kostina2019} V.~Kostina and B.~Hassibi, ``Rate-cost tradeoffs in
control,'' \emph{IEEE Trans. Autom. Control}, vol.~64, no.~11,
pp.~4525--4540, 2019.

\bibitem{tanaka2015} T.~Tanaka, K.~K.~K. Kim, P.~A. Parrilo, and S.~K.
Mitter, ``Semidefinite programming approach to Gaussian sequential
rate-distortion trade-offs,'' \emph{IEEE Trans. Autom. Control}, vol.~62,
no.~4, pp.~1896--1910, 2017.

\bibitem{attia2018} A.~Attia, A.~Alexanderian, and A.~K. Saibaba,
``Goal-oriented optimal design of experiments for large-scale Bayesian linear
inverse problems,'' \emph{Inverse Problems}, vol.~34, no.~9, 095009, 2018.

\bibitem{wu2023} K.~Wu, P.~Chen, and O.~Ghattas, ``An offline-online
decomposition method for efficient linear Bayesian goal-oriented optimal
experimental design,'' \emph{SIAM J. Sci. Comput.}, vol.~45, no.~1,
pp.~B57--B77, 2023.

\bibitem{quadgoed2025} ``Goal-oriented optimal experimental design using
quadratic approximations of nonlinear goal functionals,'' arXiv:2411.07532,
2025 (J. Sci. Comput., to appear).

\bibitem{shah2022} S.~Shah, K.~Wang, B.~Wilder, A.~Perrault, and M.~Tambe,
``Decision-focused learning without decision-making: Learning locally
optimized decision losses,'' in \emph{Advances in Neural Information
Processing Systems}, 2022.

\bibitem{shah2024} S.~Shah, B.~Wilder, A.~Perrault, and M.~Tambe, ``Leaving
the nest: Going beyond local loss functions for predict-then-optimize,'' in
\emph{Proc. AAAI Conf. Artificial Intelligence}, 2024.

\bibitem{bansal2023} D.~Bansal, R.~T.~Q. Chen, M.~Mukadam, and B.~Amos,
``TaskMet: Task-driven metric learning for model learning,'' in \emph{Advances
in Neural Information Processing Systems}, 2023.

\bibitem{elmachtoub2022} A.~N. Elmachtoub and P.~Grigas, ``Smart `predict,
then optimize','' \emph{Management Science}, vol.~68, no.~1, pp.~9--26, 2022.

\bibitem{amos2017} B.~Amos and J.~Z. Kolter, ``OptNet: Differentiable
optimization as a layer in neural networks,'' in \emph{Proc. Int. Conf.
Machine Learning}, 2017.

\bibitem{mandi2024} J.~Mandi \emph{et al.}, ``Decision-focused learning:
Foundations, state of the art, benchmark and future opportunities,'' \emph{J.
Artif. Intell. Res.}, vol.~80, pp.~1623--1701, 2024.

\bibitem{lacostejulien2011} S.~Lacoste-Julien, F.~Husz\'ar, and
Z.~Ghahramani, ``Approximate inference for the loss-calibrated Bayesian,'' in
\emph{Proc. Int. Conf. Artificial Intelligence and Statistics}, 2011.

\bibitem{zou2022} H.~Zou, C.~Zhang, S.~Lasaulce, L.~Saludjian, and H.~V.
Poor, ``Goal-oriented quantization: Analysis, design, and application to
resource allocation,'' \emph{IEEE J. Sel. Areas Commun.}, vol.~41, no.~1,
pp.~42--54, 2023.

\bibitem{shlezinger2019} N.~Shlezinger, Y.~C. Eldar, and M.~R.~D. Rodrigues,
``Hardware-limited task-based quantization,'' \emph{IEEE Trans. Signal
Process.}, vol.~67, no.~20, pp.~5223--5238, 2019.

\bibitem{zhang2018} R.~Zhang, P.~Isola, A.~A. Efros, E.~Shechtman, and
O.~Wang, ``The unreasonable effectiveness of deep features as a perceptual
metric,'' in \emph{Proc. IEEE/CVF Conf. Computer Vision and Pattern
Recognition}, 2018.

\bibitem{sun2020} Y.~Sun, Y.~Polyanskiy, and E.~Uysal, ``Sampling of the
Wiener process for remote estimation over a channel with random delay,''
\emph{IEEE Trans. Inf. Theory}, vol.~66, no.~2, pp.~1118--1135, 2020.

\bibitem{maatouk2022} A.~Maatouk, M.~Assaad, and A.~Ephremides, ``The age of
incorrect information: An enabler of semantics-empowered communication,''
\emph{IEEE Trans. Wireless Commun.}, vol.~22, no.~4, pp.~2621--2635, 2023.

\bibitem{kaymak2018} Y.~Kaymak, R.~Rojas-Cessa, J.~Feng, N.~Ansari, M.~Zhou,
and T.~Zhang, ``A survey on acquisition, tracking, and pointing mechanisms
for mobile free-space optical communications,'' \emph{IEEE Commun. Surveys
Tuts.}, vol.~20, no.~2, pp.~1104--1123, 2018.

\end{thebibliography}

\end{document}
